\begin{problem}[33届物理竞赛复赛]{62}
    两根质量均匀分布的杆 AB 和 BC，质量均为 m，长均为 l，A 端被光滑铰接到一固定点（即 AB 杆可在竖直平面内绕 A 点无摩擦转动）。开始时 C 点有外力保持两杆静止，A、C 在同一水平线 AD 上，A、B、C 三点都在同一竖直平面内， $\angle ABC = 60^{\circ}$ 。某时刻撤去外力后两杆始终在竖直平面内运动。

    \begin{subproblem}
        若两杆在 B 点固结在一起，求

        \begin{subsubproblem}
            初始时两杆的角加速度；
        \end{subsubproblem}
        \begin{subsubproblem}
            当 AB 杆运动到与水平线 AD 的夹角为 $\theta$ 时，AB 杆绕 A 点转动的角速度。
        \end{subsubproblem}
    \end{subproblem}
    \begin{subproblem}
        若两杆在 B 点光滑铰接在一起（即 BC 杆可在竖直平面内绕 B 点无摩擦转动），求初始时两杆的角加速度以及两杆间的相互作用力。
    \end{subproblem}
\end{problem}